% from ega4-8-fr.tex
\oldpage[IV-3]{6}
donc dans $Y$. Les résultats du présent paragraphe impliqueront que la
géométrie algébrique des $\mathcal O_y$-préschémas de présentation finie (et
la théorie des Modules de présentation finie sur ces préschémas) est
essentiellement équivalente à la géométrie algébrique des préschémas de
présentation finie sur des voisinages ouverts «~assez petits~» de $y$. Ainsi,
l'énoncé (8.10.5, (xii)) implique que si un morphisme $f:X\to Y$ est de
présentation finie, alors, pour que
$X\times_Y\operatorname{Spec}(\mathcal O_y)$ soit propre sur
$\operatorname{Spec}(\mathcal O_y)$, il faut et il suffit qu'il existe un
voisinage ouvert $U$ de $y$ dans $Y$ tel que $f^{-1}(U)$ soit propre sur $U$.

Un cas particulièrement important, et dans une certaine mesure classique,
est celui où $Y$ est intègre et où $y$ est son point générique, de sorte que
$\mathcal O_y$ n'est autre que \emph{le corps $R(Y)=K$ des fonctions
rationnelles sur $Y$}. Les résultats du présent paragraphe reviennent alors à
interpréter la géométrie algébrique sur $K$ en termes de la géométrie
algébrique au-dessus d'ouverts non vides «~assez petits~» de $Y$, c'est-à-dire,
intuitivement, en termes de «~familles~» d'objets géométriques indexées par les
points d'un tel ouvert. Ce point de vue a d'ailleurs été couramment utilisé
depuis longtemps, non seulement en Géométrie algébrique sur des corps
algébriquement clos, mais aussi dans l'étude arithmétique des variétés définies
sur un \emph{corps de nombres} $K$ (extension finie de $\mathbf Q$), en
considérant ce dernier comme le corps des fractions de son anneau des entiers
$A$ («~théorie de la \emph{réduction modulo} $\mathfrak p$~» ($\mathfrak p$
idéal premier de $A$)~; cf. (I, 3.7)). Les résultats des \S\S~8 et 9
fournissent donc entre autres des fondements du langage de la «~réduction
modulo $\mathfrak p$~» en arithmétique.

On notera que dans l'exemple envisagé ici, les morphismes
$S_\mu\to S_\lambda$ (pour $\lambda\leq\mu$) sont les \emph{immersions
ouvertes} canoniques $U_\mu\to U_\lambda$, et \emph{a fortiori} sont des
morphismes \emph{plats} (mais non fidèlement plats en général), ce qui
explique l'intérêt des énoncés qui font appel à une telle restriction.

\item[\emph{b)}] Supposons que les $A_\lambda$ soient des \emph{corps}, de
sorte que $A=\varinjlim A_\lambda$ est aussi un corps. Ce cas se présente
généralement lorsqu'on part de données géométriques au-dessus d'un corps
quelconque $K$, que l'on considère comme extension d'un corps $k$ (par exemple
le sous-corps premier de $K$). Il y a alors intérêt à considérer $K$ comme
limite inductive de ses sous-extensions qui sont \emph{de type fini sur $k$},
ce qui permet dans de nombreuses questions de se ramener au cas où $K$ est
une extension de type fini de $k$. Utilisant aussi la méthode esquissée dans
\emph{a)}, on peut alors se ramener généralement au cas d'un anneau de base
$A$ qui est une \emph{algèbre intègre de type fini sur $k$}.

On notera que dans cet exemple, les morphismes $S_\mu\to S_\lambda$ sont
\emph{fidèlement plats}.

\item[\emph{c)}] Supposons qu'on s'intéresse aux propriétés, locales sur
$Y$, des préschémas de présentation finie au-dessus d'un préschéma quelconque
$Y$, que l'on peut dès lors supposer affine d'anneau $A$. Il y a alors intérêt
à considérer $A$ comme limite inductive de ses sous-anneaux qui sont des
$\mathbf Z$-algèbres \emph{de type fini}, ce qui permet de ramener de
nombreuses questions au cas où $Y$ est le spectre d'une telle algèbre. C'est
là l'explicitation du «~point de vue kroneckerien~» suivant lequel la
Géométrie algébrique se réduit à la Géométrie algébrique des préschémas de
type fini sur $\mathbf Z$ (que l'on qualifie parfois de
